\section{Introduction and scientific framing}
\label{sec:intro}

\subsection{The problem: heterogeneity inside the diagnostic label}
Bipolar disorder (BP), schizophrenia (SZ), and major depressive disorder (DR) are defined by polythetic,
consensus criteria. Two patients with the same diagnosis can share few symptoms; one patient's biology,
cognition, and course can resemble a patient with a different diagnosis more than a same-diagnosis peer. For
treatment selection and prognosis this categorical resolution is too coarse: the label predicts the
\emph{average} patient, not the patient in front of the clinician. The response, developed over the last decade
in the NIMH Research Domain Criteria (RDoC) and the Hierarchical Taxonomy of Psychopathology
(HiTOP)~\cite{insel2010rdoc,kotov2017hitop}, is to describe psychopathology along \emph{continuous dimensions}
that cut across diagnostic boundaries and units of analysis (behaviour, cognition, physiology).

FACE-ATLAS operationalizes that idea on the FACE cohorts. It separates two questions that are routinely
conflated:
\begin{itemize}\setlength{\itemsep}{1pt}
  \item \textbf{Measurement} --- where does each patient sit on a small set of continuous, diagnosis-agnostic
  axes of clinical and biological variation? This \emph{builds the map}.
  \item \textbf{Stratification} --- which regions of that map define decision-relevant subgroups? This
  \emph{draws validated territories on the map}.
\end{itemize}
The two must not be collapsed. Clustering raw instruments and naming the clusters is a \emph{theory-imposed
scoring system}, not discovery; it bakes the answer into the feature choice. FACE-ATLAS instead estimates the
measurement model first and lets stratification act on the estimated coordinates later.

\subsection{How the report is organized: one map, five questions}
The work is a single argument, built as five questions that depend on one another --- each can be asked
honestly only once the previous one is answered, so the order is logical, not administrative. The conceptual
pipeline runs in four layers (cohorts $\to$ dimensions $\to$ strata $\to$ prognosis/treatment); the five
questions below walk through it, and each is a chapter of this report.

\programstrip{0}

\begin{enumerate}\setlength{\itemsep}{1pt}
  \item[\textbf{Exist?}] \textbf{Transdiagnostic dimensions} --- the measurement map, built on the baseline
  visit (Chapters~\ref{sec:data}--\ref{sec:scoring}).
  \item[\textbf{Organize?}] \textbf{Validated strata} --- the decision regions those coordinates support
  (Chapter~\ref{sec:strata}).
  \item[\textbf{Persist?}] \textbf{Temporal coherence} --- whether the coordinates hold across the follow-up
  visits (Chapter~\ref{sec:temporal}).
  \item[\textbf{Predict?}] \textbf{Prognosis} --- incremental predictive value beyond diagnosis and severity
  (Chapter~\ref{sec:prognosis}).
  \item[\textbf{Guide treatment?}] \textbf{Treatment} --- whether a stratum moderates treatment response
  (Chapter~\ref{sec:treatment}).
\end{enumerate}

\subsection{The method in one paragraph: hybrid discovery}
The ten candidate dimensions are a \emph{soft prior ontology}, not fixed outputs. Theory proposes a candidate
structure; the data estimate the actual factor structure, loadings, cross-loadings, hierarchy, and validity. A
single global, missingness-aware Bayesian sparse bifactor / ESEM model with mixed likelihoods is fit to the
harmonized baseline, and the data are free to \emph{confirm, split, merge, downgrade (proxy), reject, or declare
not-testable} any candidate. The soft priors are the formal mechanism of ``theory suggests, data decides'':
most cross-loadings are shrunk toward zero, but the likelihood can pull any of them away if the data demand it.

\begin{definitionbox}[Three objects kept distinct throughout]
\textbf{Candidate construct} --- a clinical idea, encoded as a soft prior over loadings.\quad
\textbf{Empirical dimension} --- a stable, data-supported axis of covariance.\quad
\textbf{Patient stratum} --- a recurring, decision-relevant region of the map (the subject of the stratification chapter).
A candidate is a hypothesis; a dimension is a result; a stratum is an action boundary.
\end{definitionbox}

Each candidate runs the same gauntlet (Figure~\ref{fig:discovery}): it is proposed as a soft prior, enters one
global fit, and receives a verdict. On the FACE data this loop confirmed nine dimensions (one a \emph{split} of
the biology candidate, one a \emph{proxy} relabelling), rejected one (anhedonia), and declared three
\code{not_testable} for want of indicators --- the map is what survived adversarial checking, not what was
assumed.

\begin{figure}[h]\centering
\discoveryflow
\caption{Hybrid discovery as a verdict pipeline. Theory proposes; one global fit disposes. Each candidate is
\emph{confirmed}, \emph{split}, relabelled a \emph{proxy}, \emph{rejected}, or declared \emph{not-testable},
with the FACE outcomes shown. This is the formal sense in which the data are allowed to overturn the theory.}
\label{fig:discovery}
\end{figure}

\subsection{Load-bearing invariants}
Three commitments are enforced at every layer of the system, from the data loader to the likelihood. They are
the difference between a defensible map and an artefactual one.

\begin{invariantbox}
\textbf{(I1) No naive imputation.} Structure is estimated from each patient's \emph{observed} cells via an
observed-data likelihood; no completed patient$\times$variable matrix is ever created. Naive imputation
manufactures covariance and fabricates strata when missingness is cohort-, site-, or severity-dependent --- as
it is in FACE.\\[3pt]
\textbf{(I2) Diagnosis is metadata.} BP/SZ/DR labels are never indicators. They enter only as covariates, as a
measurement-invariance grouping variable, and as a post-hoc validation layer.\\[3pt]
\textbf{(I3) Baseline defines; follow-up validates.} Dimensions are estimated on the baseline visit V0 only;
visits V1--V4 are reserved for temporal coherence and prognosis, used \emph{after} the map is fixed.
\end{invariantbox}

A fourth, structural commitment supports the first three: there is \emph{one} global factor structure on one
harmonized matrix. An instrument absent for a cohort is simply a missing cell under the observed-data
likelihood --- there are no cohort-specific ``modules.'' Coverage is emergent and flagged per patient, never a
structural sub-model. And crucially, \emph{a dimension requires indicators}: the model can discover a
statistical dimension but cannot discover a construct that is not measured; a candidate with no valid indicators
is reported \code{not_testable}, never back-filled with an invented proxy.

\subsection{What this report contributes, and how to read it}
The claim, in one sentence: \emph{the FACE cohorts hide a real, stable, transdiagnostic biological map --- a
graded continuum, not a set of biotypes --- on which biological load is nearly independent of clinical
severity; the map forecasts two-year functioning at the group level beyond diagnosis and severity, yet, on
observational data, does not reliably guide treatment.} Scientific validity is demonstrated; strong
individual-level clinical utility is not --- and the modest and negative results are reported as the
contribution, not buried beneath the positive ones. This document is the methods-of-record and the as-built
record for the whole program, written at a level of detail beyond a journal article so that it can both seed
the eventual manuscript and serve as the project's standing technical reference. Chapter~\ref{sec:data}
describes the data foundation and the no-imputation contract.
Chapter~\ref{sec:model} specifies the measurement model: the candidate ontology, the soft priors, and the
bifactor/ESEM generative model with mixed likelihoods. Chapter~\ref{sec:estimation} develops the estimation
strategy --- the staged continuation, the Gaussian-block marginalization and its identity with FIML, the
low-rank Woodbury likelihood, and the compute strategy that makes full-$N$ inference tractable.
Chapter~\ref{sec:results} reports the as-built empirical map (\textbf{nine dimensions}) and the candidate
adjudication. Chapter~\ref{sec:validation} presents the full hardening chain --- estimator/prior-robustness
confirmation, multi-seed certification, resample robustness, cross-cohort measurement invariance, and
absolute-fit posterior-predictive checks across both likelihood blocks. Chapter~\ref{sec:scoring} projects the
map to per-patient coordinates with reliability flags. Chapter~\ref{sec:strata} asks whether those
coordinates harbour discrete biotypes (they do not: a continuum) and develops the eight extreme phenotypes,
the severity spine, and the DSM-5 head-to-head that describe it. Chapter~\ref{sec:temporal} scores the fixed
map and strata onto the follow-up visits and tests whether they \emph{persist}; Chapter~\ref{sec:prognosis}
tests whether they \emph{predict} two-year outcomes beyond diagnosis and severity; and
Chapter~\ref{sec:treatment} tests whether a stratum \emph{moderates} treatment response.
Chapter~\ref{sec:engineering} documents the software: repository structure, the config-first design, the test
suite, and the two engineering bugs found and fixed \emph{with measurement}. Chapter~\ref{sec:roadmap}
collects the future studies the work points to, and Chapter~\ref{sec:limitations} states the limitations and
open problems. Appendices collect notation,
the prior and empirical atlases, the acceptance-gate definitions, the full factor-correlation matrix, and a
glossary of instruments.

\begin{methodsbox}[Reading conventions]
Coloured callouts flag \statuspill{facekr}{key results}, \statuspill{facedef}{definitions},
\statuspill{faceinv}{invariants}, \statuspill{facemeth}{methods notes}, \statuspill{facecav}{caveats}, and
\statuspill{faceopen}{open problems}. Status pills mark each finding's evidential standing: \certified\
(passes the acceptance gate at full $N$), \provisional\ (a checkpoint or sub-sample read), \rejected\ (a
candidate the data declined). Code identifiers and file paths are set in \code{monospace}.
\end{methodsbox}
